\section{Rangkuman}
\begin{itemize}
    \item Otentikasi entitas membuktikan identitas pihak pada saat sesi dibangun (dengan \textit{nonce} agar tahan \textit{replay}), sedangkan otentikasi pesan membuktikan keaslian dan keutuhan tiap pesan.
    \item Tanda tangan digital mengikat identitas pengirim pada isi pesan melalui hash yang ditandatangani dengan kunci privat; verifikasinya memakai kunci publik.
    \item RSA menandatangani dengan $s = H(m)^d \bmod n$ dan diverifikasi dengan $H(m) = s^e \bmod n$; DSS memakai DSA/ECDSA pada subgrup prima dengan tanda tangan $(r,s)$.
    \item Pesan selalu di-hash lebih dulu karena hashing memampatkan pesan, mengefisienkan komputasi, dan mengikat tanda tangan pada isi pesan; keamanannya bergantung pada ketahanan kolisi ($2^{n/2}$).
    \item \textit{Non-repudiation}, yaitu jaminan pengirim tidak dapat menyangkal pesannya, hanya dapat dijamin oleh tanda tangan digital, bukan oleh MAC yang bersifat simetris.
    \item Sertifikat X.509 mengikat identitas subjek dengan kunci publiknya dan ditandatangani CA; keabsahannya diperiksa melalui rantai kepercayaan Root CA $\to$ Intermediate CA $\to$ sertifikat server.
    \item Jabat tangan TLS memakai kriptografi asimetris untuk otentikasi dan penyepakatan kunci, lalu kriptografi simetris untuk sesi; TLS 1.3 mempersingkat proses dan mewajibkan \textit{forward secrecy}.
\end{itemize}